\section{Rapport : 19/05/2026 \& 26/05/2026}
\subsection{Équations de McKean--Vlasov}

Durant ces deux réunions, nous avons étudié un exemple classique de limite de champ moyen : les équations de McKean--Vlasov. L'objectif est de comprendre comment décrire le comportement collectif d'un système composé d'un très grand nombre de particules en interaction.

On considère dans un premier temps un système déterministe de $N$ particules. Après une remise à l'échelle adaptée, la dynamique peut s'écrire sous la forme
\begin{equation}
	\frac{d}{dt} Z_i^{(N)}(t)
	=
	\frac{1}{N}
	\sum_{j=1}^{N}
	F\!\left(
	Z_i^{(N)}(t),
	Z_j^{(N)}(t)
	\right),
	\qquad i=1,\dots,N,
	\label{eq:mean_field}
\end{equation}
où $F$ désigne le noyau d'interaction entre les particules.

Lorsque le nombre de particules devient très grand, il n'est plus raisonnable d'étudier individuellement chacune des trajectoires. On introduit alors la \emph{mesure empirique} du système
\[
	\mu_t^{N}
	=
	\frac1N
	\sum_{i=1}^{N}
	\delta_{Z_i^{(N)}(t)},
\]
qui décrit la répartition des particules dans l'espace des états.

Sous des hypothèses de régularité suffisantes sur le noyau $F$, cette mesure empirique satisfait au sens des distributions l'équation de Vlasov
\begin{important}
	\begin{equation}
		\partial_t \mu_t
		+
		\nabla \cdot
		\bigl(
		\mu_t\,F[\mu_t]
		\bigr)
		=
		0,
		\label{eq:vlasov}
	\end{equation}
	où
	\[
		F[\mu](x)
		=
		\int_{\R^d}
		F(x,y)\,\mu(dy).
	\]
	On dira qu'une famille de mesures $(\mu_t)_{t\geq0}$ est solution de l'équation de Vlasov si, pour toute fonction test $\varphi\in C_c^1(\R^d)$, l'application
	\[
		t\longmapsto \int_{\R^d}\varphi(x)\,\mu_t(dx)
	\]
	est dérivable et vérifie
	\[
		\frac{d}{dt}
		\int_{\R^d}
		\varphi(x)\,\mu_t(dx)
		=
		\int_{\R^d}
		F[\mu_t](x)\cdot \nabla\varphi(x)\,
		\mu_t(dx),
	\]
\end{important}


L'idée principale est alors d'étudier le comportement de la famille
$(\mu^N)_{N\geq 1}$ lorsque $N$ tend vers l'infini. Si la mesure empirique initiale converge vers une mesure limite $\mu_0$, on cherche à montrer que
\[
	\mu_t^N
	\Longrightarrow
	\mu_t,
\]
où $(\mu_t)_{t\geq0}$ est la solution de l'équation de Vlasov associée à la condition initiale $\mu_0$.

Pour étudier cette convergence, nous avons rappelé que l'espace des mesures de probabilité est naturellement muni de la topologie de la convergence étroite. Une stratégie classique consiste à démontrer que la famille des lois des mesures empiriques est relativement compacte, puis à identifier toutes les valeurs d'adhérence possibles comme solutions de l'équation de Vlasov.

\begin{important}

	Le résultat de compacité utilisé repose sur la notion de tension.
	% \begin{prop}{}{}
	% 	Soit $E$ un espace polonais et $\aM \subset \aP(\aP(E))$.
	% 	Alors $\aM$ est tendue si et seulement si, pour tout $\varepsilon>0$,
	% 	il existe un compact $K \subset E$ tel que
	% 	\[
	% 		\sup_{\mathbb P\in\aM}
	% 		\mathbb P
	% 		\Bigl(
	% 		\{
	% 		\mu\in\aP(E)
	% 		:
	% 		\mu(K^c)\geq \varepsilon
	% 		\}
	% 		\Bigr)
	% 		\leq \varepsilon.
	% 	\]
	% \end{prop}
	Le théorème de Prokhorov permet alors de déduire la relative compacité de la famille considérée. Une fois cette étape obtenue, il reste à identifier les limites possibles et à établir l'unicité de la solution de l'équation de Vlasov afin de conclure à la convergence de l'ensemble de la suite.
\end{important}

Cette approche constitue un premier exemple de passage d'une description microscopique, donnée par un système de particules en interaction, à une description macroscopique gouvernée par une équation aux dérivées partielles.

\subsection{Équations différentielles stochastiques de McKean--Vlasov}

On considère maintenant le pendant stochastique du système précédent. Soient
$(W^i)_{1\leq i\leq N}$ des mouvements browniens indépendants à valeurs dans $\R^d$.
Le système de particules en interaction est donné par
\begin{important}
	\begin{equation}
		dX_t^{i,N}
		=
		\frac1N
		\sum_{j=1}^{N}
		b(X_t^{i,N},X_t^{j,N})\,dt
		+
		dW_t^i,
		\qquad i=1,\dots,N.
		\label{eq:systeme_particules_stochastique}
	\end{equation}
\end{important}
On suppose dans la suite que $b:\R^d\times\R^d\to\R^d$ est bornée et lipschitzienne.

Comme précédemment, on introduit la mesure empirique
\[
	\mu_t^N
	=
	\frac1N
	\sum_{j=1}^{N}
	\delta_{X_t^{j,N}}.
\]
Le système peut alors se réécrire sous la forme
\[
	dX_t^{i,N}
	=
	\left(
	\int_{\R^d}
	b(X_t^{i,N},y)\,\mu_t^N(dy)
	\right)dt
	+
	dW_t^i.
\]
Lorsque $N$ tend vers l'infini, on s'attend donc à ce que la mesure empirique $\mu_t^N$
soit remplacée par une mesure limite déterministe $\mu_t$. La particule typique limite
doit alors satisfaire l'équation de McKean--Vlasov
\begin{important}
	\begin{equation}
		dX_t
		=
		\left(
		\int_{\R^d}
		b(X_t,y)\,\mu_t(dy)
		\right)dt
		+
		dW_t,
		\qquad
		\aL(X_t)=\mu_t.
		\label{eq:mckean_vlasov_sde}
	\end{equation}
\end{important}
Cette équation est non linéaire au sens où le coefficient de dérive dépend de la loi
du processus que l'on cherche à construire.

Pour relier cette EDS à une équation sur les mesures, on applique la formule d'Itô à une
fonction test $\varphi\in C_c^2(\R^d)$. On obtient
\[
	d\varphi(X_t)
	=
	\nabla\varphi(X_t)\cdot
	\left(
	\int_{\R^d}
	b(X_t,y)\,\mu_t(dy)
	\right)dt
	+
	\frac12\Delta\varphi(X_t)\,dt
	+
	\nabla\varphi(X_t)\cdot dW_t.
\]
En prenant l'espérance, le terme martingale disparaît. Comme $\mu_t=\aL(X_t)$, il vient
\[
	\frac{d}{dt}
	\int_{\R^d}
	\varphi(x)\,\mu_t(dx)
	=
	\int_{\R^d}
	\left[
		\frac12\Delta\varphi(x)
		+
		\left(
		\int_{\R^d}
		b(x,y)\,\mu_t(dy)
		\right)\cdot\nabla\varphi(x)
		\right]
	\mu_t(dx).
\]
C'est cette identité qui définit l'équation de Fokker--Planck non linéaire au sens faible.

\begin{important}
	Formellement, elle correspond à l'EDP
	\begin{equation}
		\partial_t\mu_t
		=
		\frac12\Delta\mu_t
		-
		\nabla\cdot
		\left[
			\left(
			\int_{\R^d}
			b(\cdot,y)\,\mu_t(dy)
			\right)\mu_t
			\right].
		\label{eq:fokker_planck_mckean_vlasov}
	\end{equation}
\end{important}

Pour construire rigoureusement la solution de l'équation de McKean--Vlasov, on travaille
sur l'espace $C_T=C([0,T];\R^d)$ muni de la distance
\[
	d_T(x,y)
	=
	\sup_{t\in[0,T]}|x(t)-y(t)|\wedge 1.
\]
On introduit ensuite la distance de Wasserstein associée :
\[
	\aW_T(\mu,\nu)
	=
	\inf_{\pi\in\Pi(\mu,\nu)}
	\int_{C_T\times C_T}
	d_T(x,y)\,\pi(dx,dy).
\]

\begin{thm}{}{}
	Si $E$ est un espace polonais, alors $(\aP(E), \aW)$ est un espace polonais, complet pour la distance de Wasserstein. De plus, la topologie induite par cette métrique est plus fine que celle de la convergence étroite.
\end{thm}

Pour une mesure $\nu\in\aP(C_T)$, on définit $X^\nu$ comme la solution de l'EDS
\[
	X_t^\nu
	=
	X_0
	+
	\int_0^t
	\left(
	\int_{C_T}
	b(X_s^\nu,w_s)\,\nu(dw)
	\right)ds
	+
	W_t.
\]
On note alors $\Phi(\nu)=\aL(X^\nu)$. Une solution de McKean--Vlasov est exactement un
point fixe de l'application $\Phi$.

Le caractère lipschitzien de $b$ permet de montrer que les itérées de $\Phi$ deviennent
contractantes pour la distance $\aW_T$. Le théorème du point fixe donne donc l'existence
et l'unicité de la loi solution.

\subsection{Propagation du chaos}

La section précédente décrit la limite d'une particule typique. Il reste à comprendre en
quel sens le système complet de particules devient asymptotiquement indépendant lorsque
$N$ tend vers l'infini. C'est précisément l'objet de la propagation du chaos. On commence par rappeler la notion d'échangeabilité.

\begin{defn}{}{}
	Une famille de variables aléatoires
	$(X_1,\dots,X_N)$ à valeurs dans un espace polonais $E$ est dite échangeable si sa loi est
	invariante par permutation des coordonnées, c'est-à-dire si, pour toute permutation
	$\sigma$ de $\{1,\dots,N\}$,
	\[
		(X_1,\dots,X_N)
		\overset{(d)}=
		(X_{\sigma(1)},\dots,X_{\sigma(N)}).
	\]
\end{defn}

Dans le système de particules, l'indépendance n'est pas préservée par la dynamique, car les
particules interagissent. En revanche, si les conditions initiales sont échangeables, alors
cette propriété reste vraie à tout temps.

Soit maintenant $L^N\in\aP(E^N)$ une suite de lois symétriques, et soit $L\in\aP(E)$.
On dit que la suite $(L^N)_{N\geq1}$ est $L$-chaotique si, pour tout $k\geq1$ et toutes
fonctions tests $\varphi_1,\dots,\varphi_k\in C_b(E)$,
\[
	\int_{E^N}
	\varphi_1(x_1)\cdots\varphi_k(x_k)
	L^N(dx_1,\dots,dx_N)
	\longrightarrow
	\prod_{i=1}^{k}
	\int_E
	\varphi_i(x)\,L(dx).
\]
Autrement dit, tout nombre fini de particules devient asymptotiquement indépendant, et
chacune a pour loi limite $L$.

Cette notion est équivalente à la convergence de la mesure empirique.

\begin{thm}{}{}
	Si $(X_1^N,\dots,X_N^N)$ a pour loi $L^N$ et si
	\[
		\mu^N
		=
		\frac1N
		\sum_{i=1}^{N}
		\delta_{X_i^N},
	\]
	alors, sous l'hypothèse de symétrie, les assertions suivantes sont équivalentes :
	\begin{enumerate}
		\item la suite $(L^N)_{N\geq1}$ est $L$-chaotique ;
		\item la mesure empirique $\mu^N$ converge en loi vers la mesure déterministe $L$ ;
		\item la propriété de chaos est vérifiée au rang $k=2$.
	\end{enumerate}
\end{thm}

Cette équivalence montre que la propagation du chaos peut être vue comme une loi des
grands nombres pour des particules échangeables.

Revenons au système de particules
\[
	dX_t^{i,N}
	=
	\frac1N
	\sum_{j=1}^{N}
	b(X_t^{i,N},X_t^{j,N})\,dt
	+
	dW_t^i.
\]
On le compare au système découplé formé des particules de McKean--Vlasov
\[
	d\overline X_t^i
	=
	\left(
	\int_{\R^d}
	b(\overline X_t^i,y)\,\mu_t(dy)
	\right)dt
	+
	dW_t^i,
	\qquad
	\aL(\overline X_t^i)=\mu_t.
\]
On choisit le même brownien $W^i$ et la même condition initiale pour $X^{i,N}$ et
$\overline X^i$. Ce couplage permet de comparer directement les trajectoires.

Sous les hypothèses précédentes, il existe, pour tout $T>0$, une constante $C_T>0$ telle
que
\[
	\E\left[
	\sup_{t\in[0,T]}
	|X_t^{i,N}-\overline X_t^i|
	\right]
	\leq
	\frac{C_T}{\sqrt N},
	\qquad
	i=1,\dots,N.
\]
Ainsi, chaque particule du système en interaction est proche, lorsque $N$ est grand, d'une
copie indépendante du processus de McKean--Vlasov.

Cette estimation implique la propagation du chaos. En effet, pour un nombre fixé
$k\geq1$, on obtient
\[
	\aW_T
	\left(
	\aL(X^{1,N},\dots,X^{k,N}),
	\aL(\overline X^1,\dots,\overline X^k)
	\right)
	\leq
	\frac{kC_T}{\sqrt N}.
\]
Or les processus $\overline X^1,\dots,\overline X^k$ sont indépendants et ont tous pour loi
$\aL(X)$. On en déduit que
\[
	\aL(X^{1,N},\dots,X^{k,N})
	\Longrightarrow
	\aL(X)^{\otimes k}.
\]
